\subsection{PRIME IDEALS}

PROPOSITION 4. Let A be a commutative ring and p an ideal of A. The following conditions are equivalent:

\begin{enumerate}
\def\labelenumi{(\alph{enumi})}
\item
  the ring A/p is an integral domain;
\item
  \(\mathbf{A} \neq \mathfrak{p}\) and, if \(x \in \mathbf{A} - \mathfrak{p}\) and \(y \in \mathbf{A} - \mathfrak{p}\) , then \(xy \in \mathbf{A} - \mathfrak{p}\) .
\item
  \(\mathfrak{p}\) is the kernel of a homomorphism of \(\mathbf{A}\) into a field.
\end{enumerate}

The implications (c) \(\Rightarrow\) (b) \(\Rightarrow\) (a) are obvious. If A/p is an integral domain, let \(f\) be the canonical injection of A/p into its field of fractions and \(g\) the canonical homomorphism of A onto A/p; then p is the kernel of \(f \circ g\) , whence the implication (a) \(\Rightarrow\) (c).

DEFINITION 3. In a commutative ring A, an ideal p satisfying the conditions of Proposition 4 is called a prime ideal.

Examples. (1) Let A be a commutative ring. If m is a maximal ideal of A, m is prime; the ring A/m is a field (no. 1, Corollary 1 to Theorem 1).

\begin{enumerate}
\def\labelenumi{(\arabic{enumi})}
\setcounter{enumi}{1}
\tightlist
\item
  If \(\mathbf{A}\) is an integral domain, the ideal \(\{0\}\) of \(\mathbf{A}\) is prime (but not maximal in general, as the example of the ring \(\mathbf{Z}\) proves).
\end{enumerate}

